\exercisesection{Exercises for § 2}

We denote by R a reduced root system in a real vector space V.

\begin{enumerate}
\def\labelenumi{\arabic{enumi})}
\tightlist
\item
  Let \(x \in V\) and let \(W(x)\) be the subgroup of \(W(R)\) consisting of the elements \(w\) such that
\end{enumerate}

\[
w (x) - x \in \mathrm{Q} (\mathrm{R}).
\]

Show that \(W(x)\) is generated by reflections. (Use the affine Weyl group of \(R^{\prime}\) .)

\begin{enumerate}
\def\labelenumi{\arabic{enumi})}
\setcounter{enumi}{1}
\item
  Assume that R is irreducible. Let \(\{\alpha_{1},\ldots,\alpha_{l}\}\) be a basis of R, let \(\tilde{\alpha}=\sum_{i}n_{i}\alpha_{i}\) be the highest root of R and let f be the connection index of R. Show that f-1 is equal to the number of indices i such that \(n_{i}=1\) .
\item
  With the notation and assumptions of no. 3, let u be an automorphism of the affine space E that permutes the hyperplanes \(L_{\alpha,k}\) ( \(\alpha \in R, k \in Z\) ). Show that u is a displacement. (If \(u_{0}\) is the linear map associated to u, show that the transpose of u leaves R stable, and hence belongs to the group A(R).) Deduce that G is the normaliser of \(W_{a}\) in the group of automorphisms of the affine space E.
\item
  Let \(C'\) be a chamber relative to \(W(R)\) in \(V^{*}\) , and let C be the alcove with vertex 0 contained in \(C'\) . Let \(S_{a}\) (resp. S) be the set of reflections in the walls of C (resp. \(C'\) ). The pairs \((W_{a}, S_{a})\) and \((W, S)\) are Coxeter systems, and \(S \subseteq S_{a}\) . Show that an element \(w \in W_{a}\) is \((S, \varnothing)\) -reduced (Chap. IV, §3, Exerc. 3) if and only if \(w(C) \subseteq C'\) .
\end{enumerate}

¶ 5) Assume that R is irreducible, and choose a chamber C of R in V; denote by B the corresponding basis of R.

\begin{enumerate}
\def\labelenumi{\alph{enumi})}
\item
  Show that the minuscule weights (§1, Exerc. 24) of R form a system of representatives in P(R) of the non-zero elements of P(R)/Q(R). (Apply to R' the corollary of Prop. 6.)
\item
  Let X be a saturated subset of Q(R) (§ 1, Exerc. 23). Assume that X is non-empty and not reduced to \{0\}. Let p be a non-zero element of X of minimal length. Show that \(p \in R\) . (In view of a), p does not satisfy condition (ii) of Exerc. 24 of § 1; hence, there exists \(\alpha \in R\) such that \(\langle p, \alpha^{-} \rangle \geqslant 2\) , so \(p - \alpha \in \mathbf{X}\) . Since the length of \(p - \alpha\) is strictly less than that of \(p\) , \(p - \alpha = 0\) so \(p \in \mathbb{R}\) .)
\item
  Let p be a dominant weight not belonging to Q(R). Show that the saturated subset S(p) generated by p (cf.~§1. Exerc. 23) contains a unique minuscule weight. (Remark that S(p) is contained in a non-trivial class mod. Q(R); conclude by applying a) and Exerc. 24 c) of §1.)
\item
  Let p be a dominant weight not belonging to Q(R). Prove that the following two properties are equivalent:
\end{enumerate}

\begin{enumerate}
\def\labelenumi{(\roman{enumi})}
\item
  \(p\) is minuscule.
\item
  There is no dominant weight \(q \neq p\) such that p - q is a linear combination of elements of B with non-negative integer coefficients.
\end{enumerate}

(If \(q\) satisfies the conditions of (v), let \(p_1\) be a minuscule weight belonging to \(S(q)\) . Then, \(p_1 \equiv p \mod Q(R)\) . \(p_1 < p\) , and \(a\) ) shows that \(p\) is not minuscule. Hence, (i) \(\Longrightarrow\) (v). Conversely, if \(p\) is not minuscule, let \(q \in S(p) - W.p\) ; transforming \(q\) by an element of W if necessary, we can assume that \(q \in \overline{C}\) ; by Exerc. 23 of §1, \(q < p\) . \(q \neq p\) , and \(q \equiv p \mod Q(R)\) . Hence, (v) \(\Longrightarrow\) (i).

